\chapter{Damped and Driven Harmonic Oscillator}

\section*{Introduction}

The damped and driven harmonic oscillator is the single most important model system in all of physics. It describes anything that oscillates, is subject to friction, and is driven by an external force---from a child on a swing to the electrical currents in a radio receiver, from the vibration of a bridge in the wind to the response of an atom to laser light. Its equation of motion, $m\ddot{x} + c\dot{x} + kx = F_0\cos(\omega t)$, captures the essential physics of resonance, energy dissipation, and frequency response that underlies virtually every oscillatory phenomenon in nature and engineering.

A mass on a spring, oscillating in a viscous fluid and pushed by an external force, is the physical realization of this equation. The spring provides the restoring force ($-kx$), the fluid provides damping ($-c\dot{x}$), and the external agent provides the driving force ($F_0\cos\omega t$). At low driving frequencies, the mass follows the force quasi-statically, oscillating in phase with the drive. As the driving frequency approaches the natural frequency $\omega_0 = \sqrt{k/m}$, the response grows dramatically---this is resonance. At resonance, the energy input from the driving force exactly balances the energy lost to damping, and the oscillation amplitude reaches its maximum. Above resonance, the mass lags behind the force by nearly $\pi$, and the amplitude decreases.


\begin{equation}
m\ddot{x} + c\dot{x} + kx = F_0\cos(\omega t)
\end{equation}

\section*{Fundamental Equations Used}

The derivation starts from Newton's second law combined with Hooke's law for the restoring force.

\section*{Assumptions}

\textbf{Assumptions:} The restoring force is linear (Hooke's law: $F_s = -kx$). The damping force is proportional to velocity (viscous damping: $F_d = -c\dot{x}$). The driving force is sinusoidal with constant amplitude and frequency. The mass $m$, damping coefficient $c$, and spring constant $k$ are time-independent.


\textbf{Limitations:} Does not describe nonlinear oscillators (where the restoring force is not proportional to displacement---e.g., large-amplitude pendulum, van der Pol oscillator). Does not capture dry (Coulomb) friction, which is constant in magnitude rather than proportional to velocity. Does not include parametric excitation (where parameters like $k$ vary with time). The transient solution decays only if $c > 0$ (underdamped or critically damped); the undamped case ($c = 0$) has persistent transients.


\section*{Mathematical Derivation}

\textbf{Step 1: Write the equation of motion for the undamped, undriven oscillator.}

For a mass $m$ on a spring with constant $k$, Newton's second law gives:
\begin{equation}
m\ddot{x} + kx = 0
\end{equation}
The natural frequency is $\omega_0 = \sqrt{k/m}$. The general solution is:
\begin{equation}
x(t) = A\cos(\omega_0 t + \phi)
\end{equation}

\textbf{Step 2: Add viscous damping.}

A damping force proportional to velocity, $F_d = -c\dot{x}$, modifies the equation to:
\begin{equation}
m\ddot{x} + c\dot{x} + kx = 0
\end{equation}
Dividing by $m$ and defining $\gamma = c/(2m)$ (the damping rate):
\begin{equation}
\ddot{x} + 2\gamma\dot{x} + \omega_0^2 x = 0
\end{equation}
The characteristic equation is $r^2 + 2\gamma r + \omega_0^2 = 0$, with roots:
\begin{equation}
r = -\gamma \pm \sqrt{\gamma^2 - \omega_0^2}
\end{equation}
For $\gamma < \omega_0$ (underdamped): $r = -\gamma \pm i\omega_d$, where $\omega_d = \sqrt{\omega_0^2 - \gamma^2}$ is the damped frequency. The solution is:
\begin{equation}
x(t) = A\,e^{-\gamma t}\cos(\omega_d t + \phi)
\end{equation}

\textbf{Step 3: Add the driving force.}

A sinusoidal driving force $F(t) = F_0\cos(\omega t)$ gives:
\begin{equation}
m\ddot{x} + c\dot{x} + kx = F_0\cos(\omega t)
\end{equation}
The general solution is $x(t) = x_h(t) + x_p(t)$, where $x_h$ is the decaying transient and $x_p$ is the steady-state response.

\textbf{Step 4: Find the steady-state solution using complex exponentials.}

Replace $\cos(\omega t)$ with $e^{i\omega t}$ and seek a particular solution $x_p(t) = \tilde{X}\,e^{i\omega t}$. Substituting:
\begin{equation}
(-m\omega^2 + ic\omega + k)\tilde{X} = F_0
\end{equation}
Solving for the complex amplitude:
\begin{equation}
\tilde{X} = \frac{F_0}{k - m\omega^2 + ic\omega}
\end{equation}
The physical solution is the real part:
\begin{equation}
x_p(t) = A(\omega)\cos(\omega t - \phi)
\end{equation}

\textbf{Step 5: Compute the amplitude and phase.}

The magnitude of $\tilde{X}$ gives the amplitude:
\begin{equation}
A(\omega) = \frac{F_0}{\sqrt{(k - m\omega^2)^2 + (c\omega)^2}}
\end{equation}
The phase lag is:
\begin{equation}
\phi(\omega) = \arctan\frac{c\omega}{k - m\omega^2}
\end{equation}

\textbf{Step 6: Find the resonance frequency.}

To find the maximum amplitude, differentiate $A(\omega)$ with respect to $\omega$ and set to zero:
\begin{equation}
\frac{d}{d\omega}\left[(k - m\omega^2)^2 + c^2\omega^2\right] = 0
\end{equation}
Expanding and differentiating:
\begin{equation}
-4m\omega(k - m\omega^2) + 2c^2\omega = 0
\end{equation}
Solving for $\omega$ (excluding $\omega = 0$):
\begin{equation}
\omega_{\text{res}} = \sqrt{\omega_0^2 - \frac{c^2}{2m^2}} = \omega_0\sqrt{1 - \frac{1}{2Q^2}}
\end{equation}
where $Q = m\omega_0/c$ is the quality factor. For low damping ($Q \gg 1$): $\omega_{\text{res}} \approx \omega_0$.

\textbf{Step 7: Evaluate the amplitude at resonance.}

At $\omega = \omega_{\text{res}}$, the maximum amplitude is:
\begin{equation}
A_{\text{max}} = \frac{F_0/c\omega_{\text{res}}}{\sqrt{1 - 1/(4Q^2)}} \approx \frac{F_0 Q}{k} \qquad (Q \gg 1)
\end{equation}
At $\omega = \omega_0$ exactly (undamped resonance frequency):
\begin{equation}
A(\omega_0) = \frac{F_0}{c\omega_0} = \frac{QF_0}{k}
\end{equation}
The phase at $\omega = \omega_0$ is $\phi = \pi/2$---the displacement lags the force by $90^\circ$, placing the velocity in phase with the force for maximum power transfer.

\section*{Characteristics of the Equation}

The equation has three regimes depending on the damping ratio $\zeta = c/(2\sqrt{mk})$: \textit{underdamped} ($\zeta < 1$): the system oscillates with decaying amplitude at frequency $\omega_d = \omega_0\sqrt{1-\zeta^2}$; \textit{critically damped} ($\zeta = 1$): the system returns to equilibrium as fast as possible without oscillating; \textit{overdamped} ($\zeta > 1$): the system returns to equilibrium sluggishly without oscillating. The quality factor $Q = \omega_0 m/c = 1/(2\zeta)$ measures the sharpness of the resonance: high $Q$ means narrow resonance and low energy loss per cycle. The steady-state amplitude at the driving frequency is $A(\omega) = F_0/\sqrt{(k - m\omega^2)^2 + (c\omega)^2}$, with a phase lag $\phi = \arctan[c\omega/(k - m\omega^2)]$.
\section*{Solution Methods}

The general solution is the sum of the transient (homogeneous) solution and the steady-state (particular) solution. The transient is $x_h(t) = Ae^{-\zeta\omega_0 t}\cos(\omega_d t - \phi_0)$ (underdamped), which decays on a timescale $\tau = 1/(\zeta\omega_0)$. The steady-state solution is $x_p(t) = A(\omega)\cos(\omega t - \phi)$, found using the method of undetermined coefficients or complex exponentials ($F_0 e^{i\omega t}$). The full solution $x(t) = x_h(t) + x_p(t)$ shows the system evolving from initial conditions toward the steady state over several decay timescales. For design purposes, the resonance curve $A(\omega)$ and the phase curve $\phi(\omega)$ are the key outputs.
\section*{Verifications}

The harmonic oscillator model is verified every day in physics and engineering labs. A driven pendulum or mass-spring system, with measured amplitude and phase as functions of driving frequency, produces the predicted Lorentzian resonance curve and the $90^\circ$ phase shift at resonance. The quality factor can be measured by the ring-down time ($Q = \pi f_0\tau$) or by the bandwidth of the resonance ($\Delta f = f_0/Q$). RLC circuits show identical resonance behavior, measured with an oscilloscope or network analyzer. The model's predictions are accurate to the extent that the linearity and time-invariance assumptions hold.
\section*{Applications}

Seismographs (the driven oscillator is the ground motion; the mass's response records the earthquake signal), RLC electrical circuits (the exact mathematical analogue: charge replaces displacement, current replaces velocity, $1/C$ replaces $k$, $R$ replaces $c$, and $L$ replaces $m$), suspension bridges and skyscrapers (design must avoid resonance with wind or earthquake frequencies---the Tacoma Narrows collapse of 1940 was a resonance catastrophe), musical instruments (the driven string or air column resonates at specific frequencies), atomic spectroscopy (atoms as harmonic oscillators driven by electromagnetic radiation, with natural linewidth determined by $Q$), and laser physics (optical cavities are resonators with very high $Q$, selecting a narrow band of frequencies).
\section*{What Richard Feynman Would Have Asked}

\textit{``Why does a driven oscillator at resonance oscillate with a $90^\circ$ phase lag relative to the driving force?''}

Feynman would insist on understanding this physically, not mathematically. At resonance, the oscillator's natural tendency is to oscillate at $\omega_0$, and the driving force is also at $\omega_0$. If the displacement were in phase with the force ($\phi = 0$), the force would push the oscillator when it is at maximum displacement (velocity is zero), doing no work on average. If the displacement lagged by $180^\circ$ ($\phi = \pi$), the force would push against the motion, extracting energy. The $90^\circ$ lag puts the velocity in phase with the force: the force pushes when the oscillator is moving fastest, doing maximum work each cycle. Nature ``chooses'' this phase because it is the only one where the energy input exactly balances the energy dissipation at the resonance amplitude. The $90^\circ$ phase shift is the signature that the oscillator is extracting maximum energy from the drive.

\textit{``What is the connection between the quality factor $Q$ and the energy stored in the oscillator?''}

$Q$ has a beautiful physical interpretation: $Q = 2\pi \times (\text{energy stored})/(\text{energy lost per cycle})$. A high-$Q$ oscillator retains its energy for many cycles; a low-$Q$ oscillator loses it quickly. The energy decays as $e^{-t/\tau}$ where $\tau = 2Q/\omega_0$, so the energy drops to $1/e$ of its initial value in $Q/\pi$ cycles. For a quartz crystal oscillator ($Q \sim 10^6$), the energy persists for hundreds of thousands of cycles after the drive is removed---this is why quartz is used as a frequency standard in watches. For a bell ($Q \sim 10^3$), the sound decays over a few hundred cycles. For a bridge ($Q \sim 10$--$100$), energy from a wind gust is dissipated in a few oscillations---unless the wind continues to pump energy in.

\textit{``What happens to the oscillator if I drive it at a frequency far above resonance?''}

Far above resonance ($\omega \gg \omega_0$), the mass cannot respond fast enough to follow the driving force. The inertia term $m\ddot{x}$ dominates, and the amplitude drops as $A \approx F_0/(m\omega^2)$, independent of the spring constant and damping. The phase lag approaches $180^\circ$: the mass moves in the opposite direction to the force. Physically, the force is pushing the mass back and forth so rapidly that by the time the mass responds, the force has already reversed---like trying to push a child on a swing by alternating pushes and pulls at twice the natural frequency, always pushing the wrong way. This is why high-frequency vibrations are attenuated by mechanical filters: the mass simply cannot follow.

\textit{``Can the transient solution ever matter in practice?''}

Yes---and Feynman would give concrete examples. When an earthquake strikes a building, the first few oscillations are dominated by the transient response (how the building responds to the initial impulse), not the steady-state response (which would be relevant if the earthquake were a sustained sinusoidal drive of fixed frequency). The transient determines whether the building survives the initial jolt. In electrical engineering, the transient of an RLC circuit determines the settling time of a filter or the startup behavior of an oscillator. In quantum mechanics, the transient response of a two-level atom to a sudden laser pulse (the Rabi oscillation transient) is the basis of atomic clocks and quantum computing operations. The steady state is what you get after waiting ``long enough,'' but ``long enough'' may be too long in practice.

\textit{``Is there a quantum version of the damped harmonic oscillator?''}

Yes, and it is surprisingly rich. A quantum harmonic oscillator coupled to a thermal bath (the quantum analogue of viscous damping) is called the quantum damped harmonic oscillator, or the Caldeira--Leggett model. The damping is modeled as coupling to a continuum of bath oscillators (representing the thermal environment). The quantum version shows phenomena absent in the classical case: zero-point fluctuations persist even at absolute zero (the oscillator has energy $\hbar\omega/2$ even with no driving), the resonance linewidth is broadened by quantum fluctuations (spontaneous emission), and at very low temperatures, the damping becomes non-Markovian (the bath ``remembers'' its past interactions). The quantum oscillator is the foundation of quantum optics (the Jaynes--Cummings model), solid-state physics (phonons), and quantum field theory (each mode of the electromagnetic field is a quantum harmonic oscillator).

\bigskip
\hrule
\bigskip
%=============================================================================
\section*{FAQ}

\textit{Q: What exactly happens at resonance?}
A: At resonance ($\omega = \omega_0$ for low damping), the amplitude reaches its maximum value $A_{\text{max}} = F_0/(c\omega_0) = QF_0/k$. The phase lag is exactly $90^\circ$: the driving force is in phase with the velocity, not the displacement. This means the force does maximum work on the system each cycle (since power = force $\times$ velocity), pumping in energy at the maximum rate. The energy dissipated by damping exactly equals the energy input, so the amplitude remains constant.

\textit{Q: Why did the Tacoma Narrows Bridge collapse?}
A: The bridge's natural frequency of torsional oscillation was excited by wind-induced vortex shedding (the Von K\'{a}rm\'{a}n vortex street). As the oscillations grew, nonlinear effects coupled the torsional mode to other modes, and the oscillations became violent enough to tear the structure apart. The collapse is often cited as a ``resonance'' disaster, though the actual mechanism was more nuanced---it involved self-excited oscillation (aeroelastic flutter) rather than simple linear resonance. The key lesson is the same: when energy is pumped into a structure at or near its natural frequency, the response can be catastrophic.

\textit{Q: How does the $Q$ factor relate to bandwidth?}
A: The quality factor $Q = f_0/\Delta f$, where $\Delta f$ is the full width at half-maximum (FWHM) of the resonance curve. A high-$Q$ oscillator has a narrow resonance (sharp frequency selectivity); a low-$Q$ oscillator has a broad resonance. In a radio receiver, a high-$Q$ LC circuit selects a narrow band of frequencies (the station you want) while rejecting others. The trade-off: high $Q$ means the oscillator takes longer to start up and longer to decay after the drive is removed ($\tau = Q/(\pi f_0)$).
\section*{Important Bibliography}
\begin{enumerate}
\item Griffiths, D.J., \textit{Introduction to Electrodynamics}, 4th ed. (Cambridge University Press, 2017), Chapter 9.
\item Morin, D., \textit{Introduction to Classical Mechanics with Problems and Solutions}, 2nd ed. (Cambridge University Press, 2008), Chapter 12.
\item Feynman, R.P., Leighton, R.B., and Sands, M., \textit{The Feynman Lectures on Physics}, Vol. I (Pearson, 2011), Chapters 21--23.
\item Strogatz, S.H., \textit{Nonlinear Dynamics and Chaos}, 3rd ed. (CRC Press, 2024), Chapter 5.
\item Cromer, A., \textit{Theoretical Mechanics} (Brooks/Cole, 1994), Chapter 15.
\end{enumerate}

